\section{Estimators}
\label{sec:estimators}

EconIRL implements twelve estimators organized into three families.
The classical structural estimators (NFXP, CCP, MPEC) solve the forward problem exactly on tabular state spaces.
The neural structural estimators (NNES, TD-CCP, SEES) approximate the value function to scale beyond tabular.
The inverse reinforcement learning estimators (MCE-IRL, AIRL, GLADIUS, f-IRL, IQ-Learn) recover rewards from demonstrations, with behavioral cloning as a non-structural baseline.
Table~\ref{tab:estimators} summarizes the capability of each method.

\begin{table}[t]
\centering
\caption{Estimator summary. ``Direction'' indicates whether the estimator solves the forward problem (parameters to policy) or the inverse problem (demonstrations to reward). ``Needs $P$'' indicates whether known transition matrices are required. ``SEs'' indicates whether valid standard errors are available.}
\label{tab:estimators}
\small
\begin{tabular}{lllcccc}
\toprule
Estimator & Reference & Direction & Needs $P$ & SEs & Scales & Transfer \\
\midrule
NFXP-NK    & \citet{rust1987, iskhakov2016}        & Forward & Yes & Analytical & No  & No \\
CCP-NPL    & \citet{hotzMiller1993, aguirregabiriaMira2002} & Forward & Yes & Hessian    & No  & No \\
MPEC       & \citet{suJudd2012}                     & Forward & Yes & Hessian    & No  & No \\
\midrule
NNES       & \citet{nguyen2025}                     & Forward & Yes & Valid      & Yes & No \\
TD-CCP     & \citet{adusumilliEckardt2025}          & Forward & No  & Robust     & Yes & No \\
SEES       & \citet{luoSang2024}                    & Forward & Yes & Marginal   & Yes & No \\
\midrule
MCE-IRL    & \citet{ziebart2010}                    & Inverse & Yes & Bootstrap  & No  & No \\
AIRL       & \citet{fu2018}                         & Inverse & No  & No         & No  & Yes \\
GLADIUS    & \citet{kang2025}                       & Inverse & Yes & Projected  & Yes & No \\
f-IRL      & \citet{ni2022}                         & Inverse & Yes & No         & No  & No \\
IQ-Learn   & \citet{garg2021}                       & Inverse & Yes & No         & No  & No \\
BC         & \citet{pomerleau1991}                  & Imitation & No & No        & No  & No \\
\bottomrule
\end{tabular}
\end{table}


\subsection{Classical Structural Estimators}

\paragraph{NFXP-NK.}
The nested fixed point algorithm maximizes the exact log-likelihood $\ell(\theta) = \sum_i \log \pi(a_i | s_i; \theta)$ where choice probabilities come from solving the Bellman equation to machine precision at each outer step.
The inner loop uses the SA-then-NK polyalgorithm of \citet{iskhakov2016}, which starts with successive approximation for global convergence and switches to Newton-Kantorovich near the fixed point for quadratic convergence.
Gradients of the log-likelihood with respect to $\theta$ are computed analytically via the implicit function theorem, differentiating through the Bellman fixed point without backpropagation or finite differences.
The outer loop uses BHHH optimization, which guarantees a positive-definite Hessian approximation from the outer product of per-observation score vectors.
Algorithm~\ref{alg:nfxp} shows the main loop.

\begin{algorithm}[t]
\DontPrintSemicolon
\caption{NFXP-NK \citep{rust1987, iskhakov2016}}
\label{alg:nfxp}
\KwIn{Panel $\mathcal{D}$, features $\phi$, transitions $P$, discount $\beta$}
Initialize $\theta_0$\;
\For{$k = 0, 1, 2, \ldots$}{
  \textbf{Inner loop:} Solve $V = T(V; \theta_k)$ via SA then NK\;
  Compute $Q(s,a) = u(s,a;\theta_k) + \beta \sum_{s'} P(s'|s,a)\, V(s')$\;
  Compute $\pi(a|s) = \mathrm{softmax}(Q/\sigma)$\;
  Compute $\nabla_\theta \ell$ via implicit differentiation\;
  Update $\theta_{k+1}$ via BHHH or L-BFGS-B\;
  \lIf{$\|\nabla \ell\| < \varepsilon$}{\textbf{break}}
}
\KwOut{$\hat\theta$, $V^*$, $\pi^*$, $(-H)^{-1}$ for analytical SEs}
\end{algorithm}

\paragraph{CCP-NPL.}
The Hotz-Miller inversion lemma \citep{hotzMiller1993} shows that under additive separability and IID extreme value shocks, observed choice probabilities $\hat P(a|s)$ uniquely invert to value differences via the emax correction $e(s,a) = \gamma_{\text{Euler}} - \log \hat P(a|s)$.
This eliminates the inner Bellman loop entirely.
Value recovery reduces to a single matrix inversion $(I - \beta F_\pi)^{-1}$, which costs $O(|\mathcal{S}|^3)$ once compared to thousands of value iterations in NFXP.
\citet{aguirregabiriaMira2002} show that iterating the Hotz-Miller inversion in CCP space (nested pseudo-likelihood, or NPL) converges to the MLE when the equilibrium is Lyapunov stable.
Setting $K=1$ gives a consistent but inefficient estimator.
Setting $K=5$ to $10$ recovers full MLE efficiency.
CCP-NPL is the only method in this package that extends naturally to dynamic games where computing Nash equilibria is infeasible.

\paragraph{MPEC.}
The mathematical programming with equilibrium constraints approach of \citet{suJudd2012} treats the value function $V(s)$ as explicit decision variables alongside the structural parameters $\theta$.
The Bellman equation $V = T(V;\theta)$ becomes a set of equality constraints enforced via an augmented Lagrangian.
This eliminates the inner loop entirely and jointly optimizes $(\theta, V)$ in a single L-BFGS-B call with increasing penalty on constraint violation.
On the Rust bus benchmark at $\beta = 0.95$, MPEC achieves identical parameter estimates to NFXP while running six to nine times faster.


\subsection{Neural Structural Estimators}

\paragraph{NNES.}
The neural network estimation of structural models approach of \citet{nguyen2025} uses a two-phase procedure.
In Phase 1 a neural network $V_\phi$ is trained to approximate the NPL Bellman target using data CCPs and the Hotz-Miller emax correction.
In Phase 2 the learned $V_\phi$ is plugged into the CCP log-likelihood and the structural parameters $\theta$ are optimized via L-BFGS-B.
The key theoretical contribution is the zero Jacobian property (Proposition 1): at the true CCP, $\partial \phi_{\theta_0}(P^*, V^*)/\partial P = 0$, so neural network errors in $V$ have only second-order effects on the likelihood.
This Neyman orthogonality guarantees that $\hat\theta$ is $\sqrt{n}$-consistent even when $V_\phi$ converges at the slower rate $o_p(n^{-1/4})$, and the estimator achieves the semiparametric efficiency bound without bias correction or cross-fitting.

\paragraph{TD-CCP.}
The temporal-difference CCP estimator of \citet{adusumilliEckardt2025} is the only method in this package that does not require known or estimated transition densities.
The key insight is that the CCP value function components $h(a,x)$ and $g(a,x)$ can be characterized via temporal-difference fixed-point equations that depend only on sample expectations over observed transitions $(s, a, s')$, not on the density $p(s'|s,a)$.
Algorithm~\ref{alg:tdccp} shows the linear semi-gradient variant, which solves for $h$ via a single matrix inversion.

\begin{algorithm}[t]
\DontPrintSemicolon
\caption{TD-CCP, linear semi-gradient \citep{adusumilliEckardt2025}}
\label{alg:tdccp}
\KwIn{Panel $\mathcal{D}$ with transitions $(s_t, a_t, s_{t+1})$, features $\phi$, basis $\Psi$}
Estimate CCPs $\hat P(a|s)$ from data frequencies\;
\For{each action $a$ and feature $k$}{
  Solve $(I - \beta \hat F_\pi)\, h_{a,k} = \phi_{a,k}$ via matrix inversion\;
}
Construct pseudo-outcome $Y_i = \sum_k \theta_k\, h_k(a_i, x_i)$\;
Minimize pseudo-likelihood over $\theta$ via GMM\;
Compute robust SEs via locally debiased moment condition\;
\KwOut{$\hat\theta$, robust SEs valid at $\sqrt{n}$-rate}
\end{algorithm}

The locally robust moment condition (Theorem 5 of \citealt{adusumilliEckardt2025}) uses a backward value function $\lambda$ and debiased moment $\zeta$ to correct for first-stage estimation error, producing standard errors valid at parametric rates despite nonparametric first stages.
The per-feature $h(a,x)$ decomposition also provides interpretable diagnostics showing how each feature contributes to the value function.

\paragraph{SEES.}
The sieve estimator of \citet{luoSang2024} approximates $V(s) \approx \Psi(s)\,\alpha$ where $\Psi$ is a deterministic sieve basis of Fourier or polynomial functions.
Joint L-BFGS-B optimization over $(\theta, \alpha)$ with an $\ell^2$ penalty on $\alpha$ enforces Bellman consistency without ever solving a fixed point.
The basis matrix is deterministic and continuation values reduce to a single matrix multiply $P_a \Psi$, making the computational cost $O(\text{basis\_dim})$ independent of $|\mathcal{S}|$.
On the Rust bus benchmark, SEES achieves parameter estimates within one percent of NFXP while running fifteen to thirty-four times faster.
No neural network training is involved.


\subsection{Inverse Reinforcement Learning Estimators}

\paragraph{MCE-IRL.}
Maximum causal entropy IRL \citep{ziebart2010} is the bridge estimator connecting the econometrics and machine learning traditions.
The backward pass runs soft value iteration to compute $V$, $Q$, and $\pi$ from the current reward weights $\theta$.
The forward pass computes state visitation frequencies $D_\pi$ via the occupancy measure fixed point $(I - \beta P_\pi^\top)^{-1}\rho_0$.
The gradient is the feature matching residual $\nabla_\theta \mathcal{L} = \E_{\text{data}}[\phi] - \E_\pi[\phi]$.
Algorithm~\ref{alg:mceirl} shows the main loop.

\begin{algorithm}[t]
\DontPrintSemicolon
\caption{MCE-IRL \citep{ziebart2010}}
\label{alg:mceirl}
\KwIn{Panel $\mathcal{D}$, features $\phi$, transitions $P$, discount $\beta$}
Compute empirical features $\bar f = \frac{1}{N}\sum_i \phi(s_i, a_i)$\;
Initialize $\theta_0$\;
\For{$k = 0, 1, 2, \ldots$}{
  \textbf{Backward:} Soft VI $\to$ $V$, $Q$, $\pi$ from $\theta_k$\;
  \textbf{Forward:} $D_\pi = (I - \beta P_\pi^\top)^{-1}\,\rho_0$\;
  $\bar f_\pi = \sum_s D_\pi(s) \sum_a \pi(a|s)\,\phi(s,a)$\;
  $g = \bar f - \bar f_\pi$ \tcp*{feature matching residual}
  $\theta_{k+1} = \theta_k + \alpha\, g$\;
  \lIf{$\|g\| < \varepsilon$ \textbf{and} $\|D_\pi - D_{\text{data}}\|_\infty < \tau$}{\textbf{break}}
}
\KwOut{$\hat\theta$ (reward weights), bootstrap SEs}
\end{algorithm}

This gradient is simultaneously the MCE gradient of \citet{ziebart2010} and the score of the DDC conditional logit likelihood.
An economist sees structural logit with feature matching.
A machine learning researcher sees maximum entropy IRL.
They solve the same optimization problem.
The deep variant replaces $\theta^\top\phi(s,a)$ with a neural network $r_\psi(s,a)$ following \citet{wulfmeier2016}, gaining flexibility at the cost of interpretability and standard errors.

\paragraph{AIRL.}
Adversarial inverse reinforcement learning \citep{fu2018} is the only IRL method in this package with a proven transfer guarantee.
The discriminator takes the form $D(s,a,s') = \exp(f)/(\exp(f) + \pi(a|s))$ where $f(s,a,s') = g_\psi(s) + \beta\, h_\psi(s') - h_\psi(s)$.
This structure disentangles the state-only reward $g(s)$ from the shaping potential $h(s)$.
Algorithm~\ref{alg:airl} shows the adversarial training loop.

\begin{algorithm}[t]
\DontPrintSemicolon
\caption{AIRL \citep{fu2018}}
\label{alg:airl}
\KwIn{Expert demonstrations $\mathcal{D}_E$, transitions $P$, discount $\beta$}
Initialize reward $g_\psi$, shaping $h_\psi$, policy $\pi_\omega$\;
\For{$k = 0, 1, 2, \ldots$}{
  $f(s,a,s') = g_\psi(s) + \beta\, h_\psi(s') - h_\psi(s)$\;
  $D(s,a,s') = \exp(f) \,/\, \bigl(\exp(f) + \pi_\omega(a|s)\bigr)$\;
  Update $\psi$: $\max_\psi\; \E_{\mathcal{D}_E}[\log D] + \E_\pi[\log(1-D)]$\;
  Update $\omega$: policy improvement against $f$\;
  \lIf{converged}{\textbf{break}}
}
\KwOut{$g_\psi(s)$ (disentangled reward, transfers across environments)}
\end{algorithm}

Theorem 5.1 of \citet{fu2018} proves that the state-only reward recovered by AIRL is disentangled with respect to all dynamics, meaning it transfers across environments with different transition functions.
This is the method to use when rewards learned in simulation must be deployed under different real-world dynamics.

\paragraph{GLADIUS.}
The ERM-IRL estimator of \citet{kang2025} parameterizes $Q(s,a)$ and $\zeta(s,a) = \E[V(s')|s,a]$ with neural networks and trains them with alternating mini-batch SGD.
Even-numbered batches update $\zeta$ via MSE against $\log\sum\exp(Q)$.
Odd-numbered batches update $Q$ via the negative log-likelihood.
After convergence, structural parameters are recovered via action-difference projection $\hat\theta = (\Delta\Phi^\top\Delta\Phi)^{-1}\Delta\Phi^\top\Delta Q$ where $\Delta$ denotes differences from a reference action, eliminating the unidentified additive constant.
A known limitation is that $Q$ is identified only up to a state-dependent constant $c(s)$ that leaks through asymmetric transitions, producing systematic bias in the operating cost parameter on the Rust bus benchmark.

\paragraph{f-IRL.}
The f-IRL estimator of \citet{ni2022} minimizes the $f$-divergence $D_f(\rho_E \| \rho_\pi)$ between expert and policy state-action occupancy measures over a tabular reward function.
The choice of $f$ provides a menu of robustness properties.
KL divergence gives maximum likelihood equivalence.
Chi-squared divergence gives variance sensitivity.
Total variation gives robustness to outlier demonstrations.
No features need to be specified.

\paragraph{IQ-Learn.}
The IQ-Learn estimator of \citet{garg2021} collapses the min-max IRL problem to a concave maximization over $Q$ alone, since the optimal policy is deterministically the softmax of $Q$.
Reward is recovered via the inverse Bellman operator $r(s,a) = Q(s,a) - \beta\sum_{s'} V(s')\,p(s'|s,a)$.
Section~\ref{sec:iq_learn} shows that in the tabular DDC setting this objective is algebraically equivalent to a penalized conditional log-likelihood.

\paragraph{BC (baseline).}
Behavioral cloning \citep{pomerleau1991} estimates the policy directly via frequency counting $\hat P(a|s) = N(s,a)/N(s)$.
No transitions, no Bellman equation, no reward function, and no value function are involved.
\citet{ross2011} show that BC error compounds quadratically under distribution shift at rate $O(T^2\varepsilon)$.
Any IRL method that cannot outperform BC has not learned from MDP structure.
BC establishes the floor for every benchmark in this paper.
